% !TEX root = ../main.tex
\chapter{Phụ lục và tài liệu tham khảo}

\section{Phụ lục A: Danh mục tệp mã nguồn chuẩn hóa}

\begin{table}[htbp]
\centering
\caption{Danh mục các module mã nguồn trong hệ thống -- Phần 1: Tác tử, công cụ và truy hồi}
\label{tab:source_files_part1}
\renewcommand{\arraystretch}{1.25}
\rowcolors{2}{tablerowalt}{white}
\begin{tabularx}{\textwidth}{L{4.8cm} Y}
\toprule
\rowcolor{tableheader}\color{white}\textbf{Tên tệp / Thư mục} & \color{white}\textbf{Mục đích và chức năng thực thi} \\
\midrule
\texttt{config.py} & Cấu hình tham số mô hình, Top-K, cặp ngưỡng $(T_{low}, T_{high})$, danh sách cho phép và đường dẫn hệ thống. \\
\texttt{llm.py} & Khởi tạo mô hình ngôn ngữ lớn đa nhà cung cấp (Groq, OpenAI, Ollama) theo Factory Pattern và mô hình nhúng BGE-M3. \\
\texttt{agent/state.py} & Định nghĩa cấu trúc trạng thái \texttt{AgentState} dạng TypedDict với các kênh tích lũy. \\
\texttt{agent/graph.py} & Thiết lập đồ thị StateGraph và rẽ nhánh ReAct (\texttt{agent}, \texttt{tools}, \texttt{cite\_validate}). \\
\texttt{agent/nodes.py} & Cài đặt các hàm xử lý logic tại từng nút của đồ thị LangGraph. \\
\texttt{agent/runtime.py} & Điểm vào duy nhất điều phối một lượt hội thoại: \texttt{prepare\_turn}, \texttt{persist\_turn}, \texttt{invoke\_crag}. \\
\texttt{agent/context.py} & Bộ quản lý ngữ cảnh lắp ghép lịch sử ngắn hạn (\texttt{messages}) và hồ sơ dài hạn (\texttt{memories}). \\
\texttt{tools/base.py} & Lớp cơ sở \texttt{BaseLegalTool} định nghĩa giao diện và lược đồ chuẩn cho công cụ tác tử. \\
\texttt{tools/registry.py} & Quản lý đăng ký tập trung và phân phối các lệnh gọi công cụ từ lõi tác tử. \\
\texttt{tools/crag\_search.py} & Công cụ tra cứu nội bộ: truy hồi lai, hợp nhất RRF, Cross-Encoder, bộ đánh giá ba nhánh và tinh lọc dải quy định. \\
\texttt{tools/web\_search.py} & Công cụ tìm kiếm web chính thống: viết lại truy vấn, lọc tên miền công quyền, làm sạch mã nguồn HTML. \\
\texttt{retrieval/dense.py} & Truy hồi vector ngữ nghĩa với ChromaDB và mô hình nhúng BGE-M3. \\
\texttt{retrieval/bm25.py} & Truy hồi từ khóa chính xác với thuật toán BM25Okapi. \\
\texttt{retrieval/fusion.py} & Hợp nhất danh sách xếp hạng ứng viên bằng giải thuật Reciprocal Rank Fusion ($k=60$). \\
\texttt{retrieval/reranker.py} & Mô hình Cross-Encoder BGE-Reranker-v2-m3 tái xếp hạng và tính điểm liên quan chuẩn hóa hàm Sigmoid. \\
\texttt{retrieval/refine.py} & Tinh chế tri thức: bóc tách đoạn trích thành các dải quy định pháp lý và lọc ngưỡng tối thiểu 0.40. \\
\bottomrule
\end{tabularx}
\end{table}

\begin{table}[htbp]
\centering
\caption{Danh mục các module mã nguồn trong hệ thống -- Phần 2: Nghiệp vụ, nạp tri thức, bộ nhớ và giao diện}
\label{tab:source_files_part2}
\renewcommand{\arraystretch}{1.25}
\rowcolors{2}{tablerowalt}{white}
\begin{tabularx}{\textwidth}{L{4.8cm} Y}
\toprule
\rowcolor{tableheader}\color{white}\textbf{Tên tệp / Thư mục} & \color{white}\textbf{Mục đích và chức năng thực thi} \\
\midrule
\texttt{legal/cleaner.py} & Làm sạch chuyên sâu: loại bỏ biểu mẫu chấm, số trang, khối nơi nhận, chuẩn hóa bảng mã Unicode NFC. \\
\texttt{legal/parser.py} & Phân tách cấu trúc văn bản pháp quy theo thứ bậc: Phần $\to$ Chương $\to$ Điều $\to$ Khoản $\to$ Điểm. \\
\texttt{legal/citations.py} & Bộ kiểm định trích dẫn xác định: kiểm tra sự tồn tại của \texttt{source\_id} và tính hiệu lực theo thời gian. \\
\texttt{legal/temporal.py} & Mô hình hóa lịch sử văn bản và phân giải quan hệ sửa đổi, bổ sung, thay thế. \\
\texttt{ingestion/worker.py} & Tiến trình xử lý ngầm đa luồng thực thi quy trình nạp tài liệu tự động. \\
\texttt{ingestion/chunks.py} & Thuật toán Legal-aware Chunking bảo toàn nguyên vẹn Điều luật và làm giàu tiêu đề ngữ cảnh. \\
\texttt{ingestion/indexer.py} & Quản trị đồng bộ hóa chỉ mục vector ChromaDB và chỉ mục từ khóa BM25. \\
\texttt{memory/store.py} & Quản lý cơ sở dữ liệu SQLite: phiên làm việc, nhật ký tin nhắn và thuộc tính người dùng. \\
\texttt{memory/extractor.py} & Trích xuất thuộc tính doanh nghiệp bền vững dựa trên danh mục cho phép \texttt{ALLOWED\_MEMORY\_KEYS}. \\
\texttt{eval/run\_eval.py} & Kịch bản tự động hóa đo lường kiểm chuẩn trên tập hiệu chuẩn và tập kiểm thử độc lập. \\
\texttt{eval/metrics.py} & Cài đặt công thức tính toán chỉ số truy hồi, điều hướng, RAGAS và kiểm định trích dẫn. \\
\texttt{api/main.py} & Hạ tầng FastAPI cung cấp các điểm cuối RESTful và SSE streaming thời gian thực. \\
\texttt{frontend/} & Ứng dụng Next.js 14 và TypeScript cung cấp giao diện tương tác người dùng và trang quản trị. \\
\texttt{schema.sql} & Bản đặc tả DDL của 8 bảng quan hệ chuẩn hóa trong cơ sở dữ liệu SQLite \texttt{app.db}. \\
\bottomrule
\end{tabularx}
\end{table}

\section{Phụ lục B: Phân công công việc nhóm phát triển}

\begin{table}[htbp]
\centering
\caption{Mô hình phân công trách nhiệm trong nhóm thực hiện đề tài}
\label{tab:team_roles}
\renewcommand{\arraystretch}{1.25}
\rowcolors{2}{tablerowalt}{white}
\begin{tabularx}{\textwidth}{L{4.2cm} Y}
\toprule
\rowcolor{tableheader}\color{white}\textbf{Vai trò thành viên} & \color{white}\textbf{Nội dung công việc phụ trách chính} \\
\midrule
\textbf{Thành viên 1 -- Kỹ sư Dữ liệu} & Thu thập văn bản từ cổng thông tin nhà nước; lập trình parser bóc tách điều khoản; chuẩn hóa siêu dữ liệu và quan hệ sửa đổi; thiết kế và quản trị SQLite \texttt{app.db}; xây dựng đường ống nạp dữ liệu và tích hợp Vision LLM OCR. \\
\textbf{Thành viên 2 -- Chuyên viên Truy hồi} & Triển khai kho vector ChromaDB với BGE-M3; cài đặt thuật toán BM25Okapi; lập trình giải thuật hợp nhất RRF; tích hợp mô hình xếp hạng lại Cross-Encoder BGE-Reranker-v2-m3; đo lường các chỉ số Recall@K và MRR. \\
\textbf{Thành viên 3 -- Lõi Tác tử \& CRAG} & Hiện thực bộ đánh giá truy hồi; cài đặt ba nhánh Correct, Ambiguous, Incorrect; xây dựng cơ chế tinh lọc tri thức và tìm kiếm web có kiểm soát; lập trình đồ thị tác tử LangGraph ReAct. \\
\textbf{Thành viên 4 -- Kiểm chuẩn \& Đảm bảo} & Xây dựng căn cứ sự thật cho tập hiệu chuẩn và tập kiểm thử độc lập; thực hiện quét lưới chọn cặp ngưỡng; đo lường các chỉ số RAGAS và Fallback; cấu hình giám sát LangFuse và phân tích các ca thất bại. \\
\textbf{Thành viên 5 -- Giao diện \& Bộ nhớ} & Phát triển ứng dụng Next.js 14 và API FastAPI hỗ trợ luồng SSE streaming; xây dựng hệ thống bộ nhớ phiên và hồ sơ ngữ nghĩa; thực hiện kiểm thử an toàn cách ly dữ liệu đa người dùng. \\
\bottomrule
\end{tabularx}
\end{table}

\section{Phụ lục C: Tiêu chí tự kiểm tra chất lượng trước khi nghiệm thu}

\begin{calloutbox}{Bảng kiểm tra điều kiện nghiệm thu hệ thống}
\begin{itemize}
    \item Có ít nhất một ca kiểm thử \textbf{Correct} trả lời chính xác và dẫn nguồn trích dẫn hợp lệ.
    \item Có ít nhất một ca kiểm thử \textbf{Incorrect} kích hoạt thành công tìm kiếm web có kiểm soát.
    \item Có ít nhất một ca kiểm thử \textbf{Ambiguous} hợp nhất thành công tri thức nội bộ và mở rộng.
    \item Có ít nhất một truy vấn siêu dữ liệu hoặc định vị điều khoản trực tiếp từ cơ sở dữ liệu quan hệ SQLite.
    \item Vết thực thi (\textit{execution trace}) thể hiện rõ điểm số đánh giá và nhánh rẽ được lựa chọn.
    \item Bảng thực nghiệm so sánh đầy đủ giữa Traditional RAG, RAG + Always Web và Full CRAG.
    \item Tập hiệu chuẩn được phân tách độc lập hoàn toàn khỏi tập kiểm thử độc lập.
    \item Bộ kiểm định trích dẫn chứng minh được khả năng phát hiện và chặn câu trả lời trích dẫn sai.
    \item Có phân tích ca thất bại thực tế kèm chẩn đoán nguyên nhân gốc rễ.
    \item Tuyệt đối không sử dụng dữ liệu bộ nhớ làm nguồn căn cứ pháp lý.
    \item Đạt chỉ số an toàn tuyệt đối: Cross-client Contamination = 0.
\end{itemize}
\end{calloutbox}

\section*{Tài liệu tham khảo}
\addcontentsline{toc}{section}{Tài liệu tham khảo}

\begin{enumerate}[label={[\arabic*]}]
    \item Yan, S.-Q., Gu, J.-C., Zhu, Y., \& Ling, Z.-H. (2024). \textit{Corrective Retrieval Augmented Generation}. arXiv preprint arXiv:2401.15884v3.
    \item Lewis, P., Perez, E., Piktus, A., Petroni, F., Karpukhin, V., Goyal, N., \dots \& Kiela, D. (2020). \textit{Retrieval-Augmented Generation for Knowledge-Intensive NLP Tasks}. Advances in Neural Information Processing Systems, 33, 9459--9474.
    \item Asai, A., Sewon, M., Jiang, Z., Chen, X., Zettlemoyer, L., \& Hajishirzi, H. (2024). \textit{Self-RAG: Learning to Retrieve, Generate, and Critique through Self-Reflection}. International Conference on Learning Representations (ICLR).
    \item Es, S., James, J., Espinosa-Anke, L., \& Schockaert, S. (2023). \textit{RAGAS: Automated Evaluation of Retrieval Augmented Generation}. arXiv preprint arXiv:2309.15217.
    \item Cormack, G. V., Clarke, C. L., \& Buettcher, S. (2009). \textit{Reciprocal Rank Fusion Outperforms Condorcet and Individual Rank Learning Methods}. Proceedings of the 32nd International ACM SIGIR Conference.
    \item Yao, S., Zhao, J., Yu, D., Du, N., Shafran, I., Narasimhan, K., \& Cao, Y. (2023). \textit{ReAct: Synergizing Reasoning and Acting in Language Models}. International Conference on Learning Representations (ICLR).
    \item Quốc hội Nước Cộng hòa Xã hội Chủ nghĩa Việt Nam. (2020). \textit{Luật Doanh nghiệp số 59/2020/QH14}.
    \item Quốc hội Nước Cộng hòa Xã hội Chủ nghĩa Việt Nam. (2019). \textit{Bộ luật Lao động số 45/2019/QH14}.
    \item BAAI (Beijing Academy of Artificial Intelligence). (2024). \textit{BGE-M3 and BGE-Reranker-v2-m3 Models Documentation}. Hugging Face Model Hub.
    \item LangChain Inc. (2024). \textit{LangGraph: Multi-Agent and Stateful Graph Orchestration Reference}. https://langchain-ai.github.io/langgraph/.
    \item Groq Inc. (2024). \textit{Groq LPU Inference Engine Documentation \& API Reference}. https://groq.com.
    \item OpenRouter. (2024). \textit{Unified LLM Interface \& Routing Documentation}. https://openrouter.ai.
\end{enumerate}
